\documentclass[12pt,a4paper]{article}
\usepackage[utf8]{inputenc}
\usepackage[T1]{fontenc}
\usepackage{amsmath,amsfonts,amssymb}
\usepackage{graphicx}
\usepackage{geometry}
\usepackage{booktabs}
\usepackage{array}
\usepackage{multirow}
\usepackage{float}
\usepackage[hidelinks]{hyperref}
\usepackage{cleveref}
\usepackage{caption}
\usepackage{enumitem}
\usepackage{makecell}

\geometry{margin=2.5cm}

\crefname{equation}{Eq.}{Eqs.}
\Crefname{equation}{Equation}{Equations}

\title{\textbf{Simulation and Optimisation Infrastructure:\\
Energy Market Operations in Renewable Energy Communities}}
\author{Energy Systems Analysis}
\date{\today}

\begin{document}

\maketitle
\tableofcontents
\newpage

%%=============================================================================
\section{Simulation and Optimisation Infrastructure}
\label{sec:infrastructure}
%%=============================================================================

This section describes the methodological and computational infrastructure
underlying the energy market simulations.  The model is designed to replicate
the Austrian wholesale and retail electricity market structure at
quarterly-hour resolution across a full calendar year.  Its core purpose is
to evaluate how different organisational arrangements — supplier structure,
Renewable Energy Community (REC) participation, and flexible storage — affect
settlement outcomes for a small community of nine nodes.

The infrastructure is deliberately structured as a sequential \emph{market
pipeline}: each stage of the Austrian energy market timeline is implemented as
a self-contained computation whose outputs become the inputs of the next stage.
This design preserves the causal ordering of real market operations and allows
each mechanism to be activated, deactivated, or modified independently across
scenarios.

%%-----------------------------------------------------------------------------
\subsection{Sequential Market Simulation Engine}
\label{subsec:engine}
%%-----------------------------------------------------------------------------

\subsubsection{Temporal Structure}

The simulation operates at a temporal resolution of $\Delta t = 0.25$~h
(15-minute intervals), consistent with the Austrian quarter-hour imbalance
settlement regime.  The full simulation horizon covers the calendar year 2016
($T = 35{,}136$ intervals), chosen for its complete availability of SimBench
load and generation profiles.  All energy quantities are expressed in kWh per
quarter-hour; power quantities in kW are converted via
$E = P \cdot \Delta t$.

\subsubsection{Market Timeline and Causal Ordering}

The market pipeline mirrors the chronological stages at which a Balance
Responsible Party (BRP) commits to and reconciles its energy position in the
Austrian market.  For scenarios involving only standard wholesale markets
(Scenarios A1, A2, B1, B2), the pipeline advances through five sequential
stages per settlement interval:

\begin{enumerate}[leftmargin=2.5em, label=\textbf{Stage~\arabic*.}]
  \item \textbf{Day-Ahead (DA) Market.}  The BRP submits forecast-based net
        positions to the exchange.  For each balancing group $g$, the
        aggregated load and generation forecasts yield a signed net commitment
        $Q_{g,\mathrm{DA}}^{t}$.  Purchase and sale volumes, and their
        associated financial commitments, are fixed at this stage.
        Formally, $C_{g,\mathrm{DA}}^{t}$ and $R_{g,\mathrm{DA}}^{t}$ are
        locked in before any intra-day revision is possible.

  \item \textbf{Intra-Day (ID) Market.}  Updated forecasts, available closer
        to real time, allow the BRP to revise its position.  The signed
        adjustment $\Delta Q_{g}^{t,\pm}$ is traded at the ID clearing price
        $\pi_{\mathrm{ID}}^{t}$.  This stage produces the \emph{closing
        position} $(\bar{Q}_{g}^{t,+},\, \bar{Q}_{g}^{t,-})$ that defines
        the BRP's contracted net demand for the settlement period.

  \item \textbf{REC Settlement} (Scenarios A2, B2 only at this stage).  After
        market trading closes, the REC operator applies an internal energy
        sharing rule based on \emph{actual} metered values.  The correction
        reduces participant-level import and export volumes proportionally,
        producing corrected metres $(\tilde{\ell}_{n}^{t},\,\tilde{g}_{n}^{t})$
        that are passed to the balancing stage.

  \item \textbf{Balancing Market.}  The TSO settles deviations between the
        BRP's contracted closing position and the realised net position
        computed from corrected meters.  Imbalance volumes
        $\varepsilon_{g}^{t}$ are valued at the system imbalance price
        $\pi_{\mathrm{imb}}^{t}$, creating penalties ($\Pi_{g}^{t}$) or
        rewards ($\Rho_{g}^{t}$) depending on sign.

  \item \textbf{Supplier Billing.}  The supplier settles individual member
        accounts using the corrected net exchange $\tilde{m}_{n}^{t}$:
        net imports are charged at the retail tariff
        $\pi_{\mathrm{ret},s(n)}^{t}$ and net exports are credited at the
        feed-in tariff $\pi_{\mathrm{fi},s(n)}^{t}$.
\end{enumerate}

\noindent
This structure ensures that: (i) contracted positions are fixed before
realisation; (ii) REC corrections act on actual metered values post-trading;
and (iii) balancing costs reflect only the residual deviation unexplained by
either the forecast revision or the REC correction.  The output of each stage
is stored as an interval-level array and propagated downstream without
modification, preserving the independence of each market mechanism.

\subsubsection{Feed-Forward Architecture}

The pipeline is strictly feed-forward: no stage modifies results produced by
a prior stage.  Any influence between stages passes exclusively through the
state arrays that connect them.  Specifically:

\begin{itemize}
  \item The DA stage outputs $Q_{g,\mathrm{DA}}^{t,\pm}$ and
        $C_{g,\mathrm{DA}}^{t}$, $R_{g,\mathrm{DA}}^{t}$.
  \item The ID stage reads the DA outputs and produces
        $\bar{Q}_{g}^{t,\pm}$, $C_{g,\mathrm{ID}}^{t}$, $R_{g,\mathrm{ID}}^{t}$.
  \item The REC stage reads actual meters and produces
        $\tilde{\ell}_{n}^{t}$, $\tilde{g}_{n}^{t}$.
  \item The balancing stage reads the closing position and the corrected
        meters to produce $\varepsilon_{g}^{t}$, $\Pi_{g}^{t}$, $\Rho_{g}^{t}$.
  \item The billing stage reads corrected meters and produces
        $\text{Bill}_{n}^{t}$ and supplier $R_{s}^{\mathrm{ret},t}$,
        $C_{s}^{\mathrm{fi},t}$.
\end{itemize}

This architecture makes the model \emph{modular}: any stage can be replaced
or disabled without altering the logic of the remaining stages.  In Scenarios
A1 and B1, the REC stage simply passes through uncorrected meters.

%%-----------------------------------------------------------------------------
\subsection{Forecasting and Uncertainty Modelling}
\label{subsec:forecasting}
%%-----------------------------------------------------------------------------

\subsubsection{Role of Forecasts in the Market Pipeline}

Forecast error is the primary driver of balancing cost in the model.  The
market pipeline requires two sets of forecasts:

\begin{itemize}
  \item \textbf{Day-ahead forecasts}
        $(\hat{\ell}_{n,\mathrm{DA}}^{t},\, \hat{g}_{n,\mathrm{DA}}^{t})$:
        used to submit DA commitments.
  \item \textbf{Intra-day forecasts}
        $(\hat{\ell}_{n,\mathrm{ID}}^{t},\, \hat{g}_{n,\mathrm{ID}}^{t})$:
        used to revise positions in the ID market and, in Scenario C3, to
        optimise the battery schedule.
\end{itemize}

Both forecast sets are derived from the same underlying SimBench profiles by
applying a controlled stochastic perturbation.

\subsubsection{Stochastic Perturbation Method}

Let $v_{n}^{t}$ denote the true (actual) value of a quantity for participant
$n$ at interval $t$.  The day-ahead forecast is generated as

\begin{equation}
\hat{v}_{n,\mathrm{DA}}^{t}
  = v_{n}^{t}\,\bigl(1 + \sigma_{\mathrm{DA}}\,\xi_{n}^{t}\bigr),
  \qquad \xi_{n}^{t} \sim \mathcal{N}(0,1)
\label{eq:da_forecast}
\end{equation}

where $\sigma_{\mathrm{DA}}$ is the day-ahead relative error standard
deviation.  The intra-day forecast applies a smaller perturbation reflecting
improved information:

\begin{equation}
\hat{v}_{n,\mathrm{ID}}^{t}
  = v_{n}^{t}\,\bigl(1 + \sigma_{\mathrm{ID}}\,\xi_{n}^{t}\bigr),
  \qquad \sigma_{\mathrm{ID}} < \sigma_{\mathrm{DA}}
\label{eq:id_forecast}
\end{equation}

The noise realisations $\xi_{n}^{t}$ are drawn from a fixed, seeded random
number generator, ensuring that all scenario comparisons are made under
identical stochastic conditions.  This is essential for isolating the effect
of market structure from the effect of forecast realisation.

\subsubsection{Asymmetric Accuracy Assumptions}

In practice, load forecasting is considerably more accurate than renewable
generation forecasting.  The model reflects this by applying distinct
$\sigma$ values to load and generation:

\begin{equation}
\sigma_{\mathrm{DA}}^{(\ell)} < \sigma_{\mathrm{DA}}^{(g)}, \qquad
\sigma_{\mathrm{ID}}^{(\ell)} < \sigma_{\mathrm{ID}}^{(g)}
\label{eq:asymmetric_sigma}
\end{equation}

Consequently, the dominant source of imbalance exposure in all scenarios is
PV generation uncertainty.  This is consistent with empirical evidence from
the Austrian balancing market.

\subsubsection{Error Propagation to Imbalance}

The relationship between forecast error and imbalance cost is direct.  The
imbalance of balancing group~$g$ at interval~$t$ is

\begin{equation}
\varepsilon_{g}^{t}
  = \underbrace{A_{g}^{t}}_{\text{actual net load}}
  - \underbrace{\bigl(\bar{Q}_{g}^{t,+} - \bar{Q}_{g}^{t,-}\bigr)}_{\text{contracted position}}
\label{eq:imbalance_error}
\end{equation}

Since the contracted position is based on intra-day forecasts, the imbalance
is approximately:

\begin{equation}
\varepsilon_{g}^{t} \approx
  \sum_{n \in \mathcal{N}_{g}} \bigl(\ell_{n}^{t} - \hat{\ell}_{n,\mathrm{ID}}^{t}\bigr)
  - \sum_{n \in \mathcal{P}_{g}} \bigl(g_{n}^{t} - \hat{g}_{n,\mathrm{ID}}^{t}\bigr)
\label{eq:imbalance_error_approx}
\end{equation}

(more precisely, after applying the REC correction in scenarios where the
REC is active).  The balancing settlement is therefore a direct financial
measure of total forecast error, weighted by the imbalance price.

\subsubsection{Performance Metrics}

Model accuracy is reported using the Mean Absolute Percentage Error (MAPE):

\begin{equation}
\mathrm{MAPE}_{n}
  = \frac{100\%}{T}
    \sum_{t \in \mathcal{T}}
    \left|
      \frac{v_{n}^{t} - \hat{v}_{n,\mathrm{ID}}^{t}}{v_{n}^{t}}
    \right|
\label{eq:mape_infra}
\end{equation}

and the Root Mean Squared Error (RMSE):

\begin{equation}
\mathrm{RMSE}_{n}
  = \sqrt{\frac{1}{T} \sum_{t \in \mathcal{T}}
    \bigl(v_{n}^{t} - \hat{v}_{n,\mathrm{ID}}^{t}\bigr)^{2}}
\label{eq:rmse}
\end{equation}

These metrics are reported separately for load and generation to quantify the
relative contribution of each source to overall imbalance.

%%-----------------------------------------------------------------------------
\subsection{Renewable Energy Community Settlement Modelling}
\label{subsec:rec_model}
%%-----------------------------------------------------------------------------

\subsubsection{Position within the Pipeline}

The REC settlement is applied \emph{after} the intra-day market closes and
\emph{before} the balancing settlement.  This ordering reflects the Austrian
regulatory framework, under which the REC sharing correction is applied to
metered values that have already been remunerated at retail and feed-in
tariffs, and the residual grid exchange is then reported to the BRP for
balancing purposes.  The correction is therefore invisible to the wholesale
market — it reduces only the bilateral supplier-customer settlement and the
BRP's residual imbalance exposure.

\subsubsection{Aggregation and Shared Energy}

For a REC $r$ with member set $\mathcal{M}_{r}$, the total generation and
load within the community at interval $t$ are:

\begin{equation}
G_{r}^{t} = \sum_{n \in \mathcal{P} \cap \mathcal{M}_{r}} g_{n}^{t},
\qquad
L_{r}^{t} = \sum_{n \in \mathcal{M}_{r}} \ell_{n}^{t}
\label{eq:rec_totals_infra}
\end{equation}

The quantity of energy that can be shared internally — without transiting the
public grid — is the lesser of total generation and total load:

\begin{equation}
E_{r}^{\mathrm{shared},t} = \min\!\bigl(G_{r}^{t},\; L_{r}^{t}\bigr)
\label{eq:shared_energy_infra}
\end{equation}

This rule follows the Austrian regulatory definition of collectively
self-consumed energy and is applied uniformly regardless of the
physical proximity of generators and consumers within the network.

\subsubsection{Proportional Allocation Rule}

The shared energy is distributed proportionally across generators and
consumers.  For prosumer $n \in \mathcal{P} \cap \mathcal{M}_{r}$, the
corrected generation that exits the community boundary is

\begin{equation}
\tilde{g}_{n}^{t} = g_{n}^{t}\left(1 - \frac{E_{r}^{\mathrm{shared},t}}{G_{r}^{t}}\right)
  \qquad \bigl(G_{r}^{t} > 0\bigr)
\label{eq:gen_correction_infra}
\end{equation}

and for member $n \in \mathcal{M}_{r}$, the corrected load that is drawn
from the public grid is

\begin{equation}
\tilde{\ell}_{n}^{t} = \ell_{n}^{t}\left(1 - \frac{E_{r}^{\mathrm{shared},t}}{L_{r}^{t}}\right)
  \qquad \bigl(L_{r}^{t} > 0\bigr)
\label{eq:load_correction_infra}
\end{equation}

The proportional rule ensures that no single member receives an
disproportionate benefit, and that the sum of corrections is consistent with
the total shared energy:

\begin{equation}
\sum_{n \in \mathcal{P} \cap \mathcal{M}_{r}}
  \bigl(g_{n}^{t} - \tilde{g}_{n}^{t}\bigr)
=
\sum_{n \in \mathcal{M}_{r}}
  \bigl(\ell_{n}^{t} - \tilde{\ell}_{n}^{t}\bigr)
= E_{r}^{\mathrm{shared},t}
\label{eq:rec_conservation}
\end{equation}

confirming that the correction is energy-conserving at the community boundary.

\subsubsection{Impact on Balancing Group Position}

The corrected meters alter the realised net position of each balancing group.
For group~$g$:

\begin{equation}
A_{g}^{t}
  = \sum_{n \in \mathcal{N}_{g}} \tilde{\ell}_{n}^{t}
  - \sum_{n \in \mathcal{P}_{g}} \tilde{g}_{n}^{t}
\label{eq:bg_actual_infra}
\end{equation}

This is the actual net demand that the BRP must have covered through its
combined DA/ID commitments.  In scenarios without a REC, corrected and
original meters coincide.

\subsubsection{Structural Consequence for Imbalance}

Because the REC correction is applied to actual meters but the BRP's contracted
position was derived from \emph{forecast} meters (without REC correction), a
systematic timing mismatch arises.  The BRP cannot incorporate future REC
corrections into its market commitments.  This mismatch enters the balancing
imbalance as an additional, structurally unavoidable component, which is one
of the central research questions addressed by the scenario comparisons in
\cref{sec:infrastructure}.

%%-----------------------------------------------------------------------------
\subsection{Multi-Supplier Structure Representation}
\label{subsec:multisupplier}
%%-----------------------------------------------------------------------------

\subsubsection{Motivation}

Scenarios B1 and B2 introduce a second supplier, transforming the balancing
structure from a single-BRP portfolio (Scenarios A1/A2) to a two-BRP system.
The structural change affects how generation and load are aggregated within
each balancing group, which in turn affects imbalance exposure even when total
community production and consumption remain unchanged.

\subsubsection{Balancing Group Assignment}

In the single-supplier scenarios (A1, A2), all nine community nodes belong
to a single balancing group $g^* = \mathcal{G}_{s^*}$.  Portfolio aggregation
benefits from the statistical cancellation of individual forecast errors:
positive and negative deviations at different nodes tend to offset each other,
reducing the aggregate imbalance below the sum of individual imbalances.

In the two-supplier scenarios (B1, B2), participants are partitioned into two
disjoint balancing groups $g_1 \subset \mathcal{G}_{s_1}$ and
$g_2 \subset \mathcal{G}_{s_2}$, with $g_1 \cup g_2 = \mathcal{N}$ and
$g_1 \cap g_2 = \emptyset$.  Each group forms an independent net position:

\begin{equation}
Q_{g_k,\mathrm{DA}}^{t} = G_{g_k,\mathrm{DA}}^{t} - L_{g_k,\mathrm{DA}}^{t},
\quad k = 1, 2
\label{eq:two_bg_net}
\end{equation}

The total DA commitment is the sum $Q_{g_1}^{t} + Q_{g_2}^{t}$, which
numerically equals the single-BRP commitment.  However, balancing penalties
are assessed per group and are proportional to the absolute imbalance of each
group independently:

\begin{equation}
\Pi_{\mathrm{total}}^{t}
  = \Pi_{g_1}^{t} + \Pi_{g_2}^{t}
  \geq \Pi_{g^*}^{t}
\label{eq:fragmentation_inequality}
\end{equation}

Inequality \eqref{eq:fragmentation_inequality} holds strictly whenever
$\varepsilon_{g_1}^{t}$ and $\varepsilon_{g_2}^{t}$ have opposite signs,
which occurs systematically if prosumers and consumers are not evenly
distributed across groups.  This \emph{fragmentation effect} is the primary
mechanism responsible for cost differences between A1 and B1, and between
A2 and B2.

\subsubsection{Portfolio Composition Effect}

The degree of fragmentation depends on how prosumers (with volatile generation)
are distributed between groups.  If all three prosumers belong to one group and
all six consumers to the other, generation and load deviations cannot cancel
across groups.  The model assigns participants to groups according to the
scenario configuration files, allowing the fragmentation magnitude to be
quantified empirically.

\subsubsection{Cross-Supplier REC Interaction (Scenarios B2)}

In Scenario B2 the REC is active while participants are distributed across two
suppliers.  The REC correction operates on the full community (all nine nodes)
regardless of supplier assignment.  The corrected meters are then attributed
back to the original balancing groups.

This creates a cross-supplier dependency: the REC allocation to a member of
group $g_1$ is determined by the aggregate generation and load of the entire
community, including members of group $g_2$.  Formally, for a member
$n \in \mathcal{M}_{r} \cap \mathcal{N}_{g_1}$:

\begin{equation}
\tilde{\ell}_{n}^{t} = \ell_{n}^{t}\left(1 - \frac{\min(G_{r}^{t}, L_{r}^{t})}{L_{r}^{t}}\right)
\label{eq:cross_bg_rec}
\end{equation}

where $G_{r}^{t}$ and $L_{r}^{t}$ include contributions from both $g_1$ and
$g_2$.  This cross-group dependency means that the corrected net position of
$g_1$ depends on the generation and load outcomes of $g_2$, introducing an
externality across supplier balancing accounts.

%%-----------------------------------------------------------------------------
\subsection{Cost and Settlement Aggregation}
\label{subsec:cost_aggregation}
%%-----------------------------------------------------------------------------

\subsubsection{Per-Interval Settlement Building Blocks}

At each interval $t$, four financial flows are computed for each balancing
group $g$ belonging to supplier $s$:

\begin{enumerate}[leftmargin=2.5em, label=(\roman*)]
  \item \textbf{Day-ahead cost/revenue}: purchase cost $C_{g,\mathrm{DA}}^{t}$
        and sale revenue $R_{g,\mathrm{DA}}^{t}$ at the DA price.
  \item \textbf{Intra-day adjustment}: additional purchase cost
        $C_{g,\mathrm{ID}}^{t}$ and additional sale revenue
        $R_{g,\mathrm{ID}}^{t}$ at the ID price.
  \item \textbf{Balancing penalty/reward}: $\Pi_{g}^{t}$ for net deficit,
        $\Rho_{g}^{t}$ for net surplus, at the imbalance price.
  \item \textbf{Retail/feed-in settlement}: retail revenue
        $R_{s}^{\mathrm{ret},t}$ from member import charges; feed-in cost
        $C_{s}^{\mathrm{fi},t}$ for crediting member exports.
\end{enumerate}

\subsubsection{Annual Aggregation}

The supplier's annual profit/loss is:

\begin{align}
\Pi_{s}^{\mathrm{annual}}
  &= \underbrace{\sum_{t}\sum_{g \in \mathcal{G}_{s}}S_{g,\mathrm{DA}}^{t}}_{\text{DA settlement}}
   + \underbrace{\sum_{t}\sum_{g \in \mathcal{G}_{s}}
     \bigl(R_{g,\mathrm{ID}}^{t} - C_{g,\mathrm{ID}}^{t}\bigr)}_{\text{ID adjustment}}
\notag\\
  &\quad
   + \underbrace{\sum_{t}\sum_{g \in \mathcal{G}_{s}}
     \bigl(\Rho_{g}^{t} - \Pi_{g}^{t}\bigr)}_{\text{balancing net}}
   + \underbrace{\sum_{t}
     \bigl(R_{s}^{\mathrm{ret},t} - C_{s}^{\mathrm{fi},t}\bigr)}_{\text{retail net}}
\label{eq:pl_decomposed}
\end{align}

This decomposition directly supports the scenario comparison framework:
differences between scenarios can be attributed precisely to the stage at which
they first arise.

\subsubsection{Individual Member Cost Calculation}

From the perspective of a community member, the annual cost is:

\begin{equation}
\Gamma_{n}^{\mathrm{annual}}
  = \sum_{t \in \mathcal{T}} \text{Bill}_{n}^{t}
  = \sum_{t \in \mathcal{T}}
    \Bigl[
      \max\!\bigl(\tilde{m}_{n}^{t},\;0\bigr)\,\pi_{\mathrm{ret},s(n)}^{t}
    - \max\!\bigl(-\tilde{m}_{n}^{t},\;0\bigr)\,\pi_{\mathrm{fi},s(n)}^{t}
    \Bigr]
\label{eq:member_annual_cost}
\end{equation}

The REC correction reduces $|\tilde{m}_{n}^{t}|$ relative to the uncorrected
net exchange, thereby reducing both imports charged at retail tariff and
exports credited at feed-in tariff.  The net saving to a consumer is always
positive; for a prosumer the net effect depends on the relative levels of
$\pi_{\mathrm{ret}}$ and $\pi_{\mathrm{fi}}$.

%%-----------------------------------------------------------------------------
\subsection{Model Validation and Consistency Checks}
\label{subsec:validation}
%%-----------------------------------------------------------------------------

The model incorporates a set of algebraic consistency checks that are applied
at the end of each simulated year.  These checks do not require external
benchmarks: they verify internal logical invariants that must hold regardless
of parameterisation.

\subsubsection{Community Energy Balance}

At every interval, total metered generation, total metered load, and net
grid exchange must be consistent:

\begin{equation}
\sum_{n \in \mathcal{P}} g_{n}^{t}
- \sum_{n \in \mathcal{N}} \ell_{n}^{t}
= \underbrace{X_{\mathrm{comm}}^{t} - I_{\mathrm{comm}}^{t}}_{\text{net export}}
\label{eq:energy_balance_check}
\end{equation}

where $X_{\mathrm{comm}}^{t}$ and $I_{\mathrm{comm}}^{t}$ are the aggregate
export and import at the community boundary.  Any numerically significant
violation of \cref{eq:energy_balance_check} indicates a data inconsistency
in the profile scaling.

\subsubsection{REC Conservation Check}

Equation~\cref{eq:rec_conservation} is verified numerically at every interval
to confirm that the total reduction in generation export equals the total
reduction in load import:

\begin{equation}
\left|\;
\sum_{n \in \mathcal{P} \cap \mathcal{M}_{r}}
\bigl(g_{n}^{t} - \tilde{g}_{n}^{t}\bigr)
\;-\;
\sum_{n \in \mathcal{M}_{r}}
\bigl(\ell_{n}^{t} - \tilde{\ell}_{n}^{t}\bigr)
\;\right| < \epsilon
\label{eq:rec_conservation_check}
\end{equation}

for a suitable numerical tolerance $\epsilon$.  A violation indicates an
implementation error in the proportional allocation logic.

\subsubsection{Contracted vs Realised Position Consistency}

The total contracted position across all balancing groups of a supplier must
equal the sum of the DA and ID commitments:

\begin{equation}
\sum_{g \in \mathcal{G}_{s}}
\bigl(\bar{Q}_{g}^{t,+} - \bar{Q}_{g}^{t,-}\bigr)
=
\sum_{g \in \mathcal{G}_{s}}
\bigl(Q_{g,\mathrm{ID}}^{t,+} - Q_{g,\mathrm{ID}}^{t,-}\bigr)
\label{eq:contracted_check}
\end{equation}

\subsubsection{Non-Negativity of Settlement Quantities}

By construction, purchase volumes $Q_{g}^{t,+}$, sale volumes $Q_{g}^{t,-}$,
penalties $\Pi_{g}^{t}$, rewards $\Rho_{g}^{t}$, and retail revenues
$R_{s}^{\mathrm{ret},t}$ must all be non-negative.  Negative values would
indicate a sign error in the settlement logic and are flagged as exceptions.

\subsubsection{Imbalance Decomposition Check}

The total annual balancing cost must be consistent with the sum of per-group,
per-interval balancing settlements:

\begin{equation}
C_{s}^{\mathrm{bal}}
= \sum_{t \in \mathcal{T}}\;\sum_{g \in \mathcal{G}_{s}} \Pi_{g}^{t}
\label{eq:bal_cost_check}
\end{equation}

\subsubsection{Comparative Scenario Consistency}

As a higher-level check, the interaction decomposition must be internally
consistent:

\begin{equation}
\bigl(C_{\text{B2}} - C_{\text{A1}}\bigr)
\stackrel{!}{=}
\Delta C^{\mathrm{REC}} + \Delta C^{\mathrm{comp}} + \Delta C^{\mathrm{int}}
\label{eq:interaction_check}
\end{equation}

This equality is definitional (it follows from the definitions of the three
components), but its numerical verification confirms that the cost aggregation
pipeline has been applied consistently across all four base scenarios.

%%-----------------------------------------------------------------------------
\subsection{Computational Structure}
\label{subsec:computation}
%%-----------------------------------------------------------------------------

\subsubsection{Time-Step Iteration}

For scenarios A1, A2, B1, B2 the market pipeline is implemented as a single
vectorised pass over all $T$ intervals.  Market prices (DA, ID, imbalance,
retail, feed-in) and profile values are pre-loaded into aligned arrays indexed
by $t$.  Each stage applies element-wise operations across the full time
dimension, avoiding explicit time-step loops where possible.  This design
enables efficient computation of the full annual settlement in a single
function call per stage.

\subsubsection{Modular Stage Activation}

A scenario configuration object specifies a Boolean activation flag for each
pipeline stage.  Stages not activated for a given scenario return
pass-through outputs (uncorrected meters, zero adjustments), ensuring that
the downstream pipeline receives consistently structured data regardless of
scenario.  This design allows the same pipeline code to execute all five
scenarios without conditional branching in the core settlement logic.

\subsubsection{Data Flow and Storage}

Intermediate outputs from each stage are stored as named arrays with
dimensions $(\text{node} \times T)$ or $(\text{group} \times T)$.  No
intermediate results are discarded until final annual aggregation is complete.
This retention policy supports post-hoc decomposition of any cost component
at any temporal granularity (interval, day, month, year) without re-running
the simulation.

\subsubsection{Algorithmic Complexity}

For $N = |\mathcal{N}|$ participants, $G = |\mathcal{G}|$ balancing groups,
and $T$ intervals, the computational complexity of the base pipeline (Stages
1--5 for Scenarios A1/A2/B1/B2) is

\begin{equation}
\mathcal{O}\bigl(T \cdot (N + G)\bigr)
\label{eq:complexity_base}
\end{equation}

For the case study, $N = 9$, $G \in \{1, 2\}$, and $T = 35{,}136$, making the
base computation negligible in wall-clock time.

%%=============================================================================
\section{Scenario C3: Battery Optimisation Extension}
\label{sec:scenario_c3}
%%=============================================================================

Scenario C3 augments the Scenario A2 pipeline with a centralised, MILP-based
battery scheduling stage inserted between the intra-day market and the REC
settlement.  This section describes the battery optimisation method and its
integration into the simulation infrastructure.

%%-----------------------------------------------------------------------------
\subsection{Position within the Pipeline}
\label{subsec:battery_position}
%%-----------------------------------------------------------------------------

The battery scheduling problem is solved \emph{after} the intra-day market
closes and \emph{before} the REC settlement.  The rationale is:

\begin{enumerate}[leftmargin=2.5em]
  \item The intra-day forecasts
        $(\hat{\ell}_{n,\mathrm{ID}}^{t},\, \hat{g}_{n,\mathrm{ID}}^{t})$
        represent the best available information and are used as inputs to the
        optimiser.
  \item The resulting battery schedule modifies the net generation profile
        before REC sharing is applied, so the REC settlement already accounts
        for battery output when computing shared energy.
  \item The balancing imbalance in C3 reflects the three compounded
        uncertainties: forecast error, REC timing mismatch, and battery
        schedule deviation from actuals.
\end{enumerate}

The full C3 pipeline is:

\begin{equation}
\text{DA}
\;\rightarrow\; \text{ID}
\;\rightarrow\; \underbrace{\text{Battery MILP}}_{\text{C3 only}}
\;\rightarrow\; \text{REC}
\;\rightarrow\; \text{Balancing}
\;\rightarrow\; \text{Billing}
\label{eq:c3_pipeline}
\end{equation}

%%-----------------------------------------------------------------------------
\subsection{Battery Optimisation Problem}
\label{subsec:battery_milp}
%%-----------------------------------------------------------------------------

\subsubsection{Decision Variables}

The optimiser determines, for each interval $t \in \mathcal{T}$:
charge power $P_{\mathrm{c}}^{t} \geq 0$; discharge power
$P_{\mathrm{d}}^{t} \geq 0$; state of charge $E^{t}$; binary mode
variable $u^{t} \in \{0,1\}$ ($u^{t} = 1$ for charging); grid import
$I^{t} \geq 0$; grid export $X^{t} \geq 0$.

\subsubsection{Objective Function}

The optimiser minimises the community's annual net grid cost, valued at the
supplier's retail tariff for imports and feed-in tariff for exports:

\begin{equation}
\min_{\{P_{\mathrm{c}}^{t},\,P_{\mathrm{d}}^{t},\,E^{t},\,u^{t},\,I^{t},\,X^{t}\}}
\quad
\sum_{t \in \mathcal{T}}
\Bigl[
  I^{t}\,\pi_{\mathrm{ret},s}^{t}
 -X^{t}\,\pi_{\mathrm{fi},s}^{t}
\Bigr]
\label{eq:battery_obj_c3}
\end{equation}

\subsubsection{State-of-Charge Dynamics}

The SoC evolves according to the charge/discharge balance:

\begin{equation}
E^{t} = E^{t-1} + \eta_{\mathrm{c}}\,P_{\mathrm{c}}^{t}\,\Delta t
              - \frac{P_{\mathrm{d}}^{t}\,\Delta t}{\eta_{\mathrm{d}}},
\qquad E^{0} = E_{0}
\label{eq:soc_dynamics_c3}
\end{equation}

with charge efficiency $\eta_{\mathrm{c}} = 0.95$ and discharge efficiency
$\eta_{\mathrm{d}} = 0.95$.  The asymmetric treatment of charge and discharge
efficiencies in \cref{eq:soc_dynamics_c3} means that a charge-discharge cycle
always incurs a net energy loss of approximately
$1 - \eta_{\mathrm{c}} \cdot \eta_{\mathrm{d}} = 9.75\%$.

\subsubsection{Physical and Operational Constraints}

\begin{align}
E_{\min} &\leq E^{t} \leq E_{\max}
&&\forall\, t \in \mathcal{T}
\label{eq:soc_bounds_c3}\\[4pt]
0 &\leq P_{\mathrm{c}}^{t} \leq P_{\mathrm{c}}^{\max}\,u^{t}
&&\forall\, t \in \mathcal{T}
\label{eq:charge_limit_c3}\\[4pt]
0 &\leq P_{\mathrm{d}}^{t} \leq P_{\mathrm{d}}^{\max}\,(1 - u^{t})
&&\forall\, t \in \mathcal{T}
\label{eq:discharge_limit_c3}
\end{align}

Constraints~\cref{eq:charge_limit_c3,eq:discharge_limit_c3} are the big-M
formulation of binary mutual exclusivity: the battery can charge or discharge
in any interval, but not both simultaneously.  This prevents the optimiser
from exploiting round-trip efficiency differentials spuriously.

The parameters for Scenario C3 are: $E_{\max} = 200$~kWh;
$P_{\mathrm{c}}^{\max} = P_{\mathrm{d}}^{\max} = 50$~kW;
$E_{\min} = 0.20 \times E_{\max} = 40$~kWh; $E_{0} = 0.50 \times E_{\max} = 100$~kWh.

\subsubsection{Net Power Balance}

Let $\tilde{L}^{t}$ and $\tilde{G}^{t}$ denote the community-level intra-day
load and generation forecasts.  The grid import/export decision variables are
linked to battery operation through:

\begin{equation}
I^{t} - X^{t}
  = \tilde{L}^{t} - \tilde{G}^{t}
    + P_{\mathrm{c}}^{t}\,\Delta t - P_{\mathrm{d}}^{t}\,\Delta t,
\qquad I^{t} \geq 0,\quad X^{t} \geq 0
\label{eq:grid_balance_c3}
\end{equation}

This constraint ensures that the optimised schedule is physically consistent
with the community's net load, as estimated from intra-day forecasts.

\subsubsection{Modelling Approach and Solver Class}

The problem is a Mixed-Integer Linear Programme (MILP) with $6T$ decision
variables ($T$ binary) and $5T$ constraints.  For $T = 35{,}136$ the problem
has approximately $211{,}000$ variables and $176{,}000$ constraints.  The
problem is solved as a single annual optimisation using the branch-and-bound
method, with the GLPK/CBC solver invoked via the Pyomo algebraic modelling
framework.

\subsubsection{Interaction with Forecast Updates}

The battery schedule is optimised using intra-day forecasts, not actual values.
Deviations of actual generation and load from the intra-day forecast therefore
cause the battery's realised impact on net load to differ from its scheduled
impact.  This forecast-induced deviation enters the balancing imbalance as an
additional source of error, specific to Scenario C3:

\begin{equation}
\varepsilon_{g,C3}^{t} = \varepsilon_{g,A2}^{t}
  + \underbrace{
    \bigl(P_{\mathrm{d,actual}}^{t} - P_{\mathrm{c,actual}}^{t}\bigr)
    - \bigl(P_{\mathrm{d,sched}}^{t} - P_{\mathrm{c,sched}}^{t}\bigr)
  }_{\text{battery schedule deviation}}
\label{eq:battery_imbalance}
\end{equation}

Whether the battery reduces or increases net balancing cost relative to
Scenario A2 depends on the relative magnitude of the cost reduction from
peak shaving and export arbitrage versus the additional imbalance from
schedule deviations.

%%-----------------------------------------------------------------------------
\subsection{C3 Validation Checks}
\label{subsec:c3_validation}
%%-----------------------------------------------------------------------------

In addition to the base-scenario checks of \cref{subsec:validation}, the
following invariants are verified for Scenario C3:

\begin{enumerate}[leftmargin=2.5em]
  \item \textbf{SoC feasibility}: $E_{\min} \leq E^{t} \leq E_{\max}$ for
        all $t$.
  \item \textbf{Binary exclusivity}: $P_{\mathrm{c}}^{t} \cdot P_{\mathrm{d}}^{t} = 0$
        for all $t$ (numerically: both cannot be simultaneously above a small
        tolerance).
  \item \textbf{Grid balance consistency}: at every $t$, $I^{t} - X^{t}$
        as returned by the solver must equal the right-hand side of
        \cref{eq:grid_balance_c3} to within numerical tolerance.
  \item \textbf{Objective consistency}: the total net grid cost reported by
        the solver must equal $\sum_t (I^t \pi_{\mathrm{ret}}^t - X^t \pi_{\mathrm{fi}}^t)$
        computed from the optimal variable values.
\end{enumerate}

%%=============================================================================
\section{Summary of Methodological Choices}
\label{sec:summary}
%%=============================================================================

\begin{table}[H]
\centering
\caption{Key methodological choices and their justification}
\label{tab:method_summary}
\begin{tabular}{@{} p{3.8cm} p{5.0cm} p{5.5cm} @{}}
\toprule
\makecell[l]{Design Choice} &
\makecell[l]{Implementation} &
\makecell[l]{Justification} \\
\midrule
Temporal resolution
  & 15-minute intervals ($\Delta t = 0.25$~h)
  & Matches Austrian imbalance settlement period \\[4pt]
Forecast perturbation
  & Seeded Gaussian noise, asymmetric $\sigma$ for load and generation
  & Reproducible; reflects empirical accuracy gap between load and PV \\[4pt]
REC timing
  & Post-ID, pre-balancing
  & Consistent with Austrian regulatory framework \\[4pt]
REC allocation rule
  & Proportional to metered share
  & Avoids discriminatory treatment of members \\[4pt]
BRP fragmentation
  & Independent net positions per group
  & Correctly captures loss of portfolio diversification \\[4pt]
Battery horizon
  & Full annual MILP
  & Provides perfect-forecast upper bound on battery value; avoids
    rolling-horizon truncation artefacts \\[4pt]
Battery binary constraint
  & Big-M formulation
  & Prevents simultaneous charge/discharge; standard MILP encoding \\[4pt]
Validation
  & Algebraic invariant checks post-simulation
  & Detects implementation errors without reference data \\
\bottomrule
\end{tabular}
\end{table}

\bibliography{references}
\bibliographystyle{plain}

\end{document}
